% Arabic abstract, compiled separately with LuaLaTeX and included as a page in the report.
\documentclass[12pt]{article}
\usepackage[margin=1in]{geometry}
\usepackage[bidi=basic, layout=tabular, provide=*]{babel}
\babelprovide[import, main]{arabic}
\babelfont{rm}[Renderer=HarfBuzz]{Times New Roman}
\usepackage{setspace}
\onehalfspacing
\pagestyle{empty}
\begin{document}
\begin{center}
{\Large\bfseries الملخص}
\end{center}
\vspace{6pt}

تؤدي التنبيهات المزعجة في أنظمة مراقبة السائق، سواء كانت صحيحة أم خاطئة، إلى قيام السائقين بتعطيل هذه الأنظمة أو التحايل عليها وفقدان ثقتهم بها، مما يعرّض السائقين النعسانين والمشتتين لخطر يهدد حياتهم. ويكمن أصل المشكلة في ضرورة تحسين هذه الأنظمة من حيث تجربة المستخدم مع الحفاظ على سلامة السائق. نراجع في هذا التقرير الأعمال السابقة وحدودها المعروفة، ونصمّم نظاماً من منظور شخص عامل متعب، أولويته الوصول إلى وجهته بأمان. تغذّي كاميرا ويب بثاً مباشراً لنموذج لكشف الأجسام محسَّن للعمل على جهاز \foreignlanguage{english}{Raspberry Pi 5}، وتقوم مراحل المعالجة القبلية والبعدية بتصفية الاكتشافات الخاطئة. وعند تأكيد أي اكتشاف، يصدر النظام تنبيهاً يتناسب مع مستوى الخطر، ضوئياً وصوتياً مع رسالة صوتية تشرح سببه. تُسجَّل الأحداث المؤكدة على الجهاز ثم تُرفع إلى لوحة تحليلات على الويب تقدّم للسائق نصائح من أجل قيادة أكثر أماناً، دون أن يغادر أي فيديو الجهاز. ويُقيَّم النظام باختبار معياري على أشخاص لم يظهروا في بيانات التدريب، يقيس الدقة والاستدعاء وعدد التنبيهات الخاطئة في الساعة.

في هذه المرحلة الأولى نقدّم تصميماً متكاملاً وليس منتجاً نهائياً. فقد حدّدنا ٢٦ متطلباً قابلاً للاختبار، وقارنّا خمسة بدائل للنظام بأكمله وفق معايير موزونة، واخترنا نموذج \foreignlanguage{english}{YOLO11n} على \foreignlanguage{english}{Raspberry Pi 5} مع مسرّع \foreignlanguage{english}{AI HAT+}، بعد أن قدّرنا أن المعالج المركزي وحده لا يتجاوز نحو ٧ إطارات في الثانية. وتُظهر دراساتنا الأولية الأربع على بيانات عامة أن \foreignlanguage{english}{YOLO11n} يتفوّق على النماذج الصغيرة الأخرى عند الاختبار على سائقين جدد، وأن إضافة ٤٠٪ من العينات السلبية الصعبة تخفض الاكتشافات الخاطئة بنسبة ٧١٪، وأن عتبات القرار المضبوطة على بعض السائقين لا تصلح لغيرهم، مما يبرّر المعايرة لكل سائق. وتشمل الأهداف ٢٥ إلى ٤٠ إطاراً في الثانية، وزمن استجابة أقل من ١٠٠ ميلي ثانية، واستدعاءً لا يقل عن ٠٫٩٠ لكل فئة، وأقل من ١٠ تنبيهات خاطئة في الساعة، ضمن ميزانية تبلغ ٢٧٢٥ درهماً. وقد بنينا البرمجيات كاملةً واختبرناها على حاسوب، فنجحت جميع الفحوصات الآلية البالغ عددها ١٩٤، ورصدت إعادة تشغيل تسجيل لسائق لم يتدرب عليه النموذج، من الفيديو حتى لوحة المعلومات عبر اتصال مشفّر، جميع حالات استخدام الهاتف الأربع وحالتين من أربع حالات شرب، دون أي إنذار خاطئ من الكاشف. كما سجّلنا أربع جلسات تجريبية لثلاثة من مؤلفي التقرير داخل سيارات حقيقية، نهاراً وليلاً. وتبقى نتائج \foreignlanguage{english}{Raspberry Pi}، التي ستُعرض على منصة مكتبية في المرحلة الثانية، للمرحلة التالية.
\end{document}
